\documentclass[10pt]{article}
\usepackage[margin=0.9in]{geometry}
\usepackage{amsmath,amssymb,amsthm}
\usepackage[hidelinks]{hyperref}
\newcommand{\WIPHASH}{\texttt{d32be4d}}
\newcommand{\file}[1]{{\footnotesize\texttt{\detokenize{#1}}}}
\newcommand{\Lead}{\mathrm{Lead}}
\newtheorem*{lemma1}{Lemma 1 (histories are increasing forests)}
\newtheorem*{lemma2}{Lemma 2 (Theorem W telescopes)}
\newtheorem*{lemma3}{Lemma 3 (tree--graph identity, generic weights)}
\newtheorem*{thmG}{Theorem G}
\newtheorem*{thmF}{Theorem F}

\title{The full-merge lead is a connected-graph cumulant\\(Theorems G and F)}
\author{Rick}
\date{2026-10-04}

\begin{document}
\sloppy
\maketitle
\vspace{-1.5em}
\noindent\textbf{Author:} Rick. \textbf{Date:} 2026-10-04.
\textbf{Recipient:} Clio Vega (cc Robin Langer).\\
\textbf{WIP commit:} \WIPHASH{} of \texttt{grandpa-rick/work-in-progress}. This is the latest commit touching the source
file \file{proofs/2026-10-05-day221-lead-cumulant.md}, and this PDF is a faithful transcription of that file.\\
\textbf{Registry grades:} nodes \texttt{conjG} and \texttt{conjF} are \emph{proved}. No one else has read them yet,
and \textbf{novelty has not been searched} (see \S8).

\paragraph{Provenance note.} The proof file committed in WIP \texttt{5b73f01} was accidentally truncated to about
1.2\,KB: the header, a list of scripts and \S8 survived, and \S\S1--7 were lost. This PDF is the restored version,
which I re-derived from scratch in commit \texttt{d32be4d} from the session outline and the Day~220 proofs of Theorems~B
and~W. That makes it a second derivation, not the original text. The text of the Day~220 Theorem~W cold re-check
(old \S6) was also lost and has \emph{not} been restored. Theorem~W is graded proved on its own Day~220 proof
(\file{proofs/2026-10-03-day220-s1-carre-du-champ.md} \S5b), but its independent re-read is still owed.

\section*{0. Setup (from Day 220 \S4)}
Let $\lambda=(\lambda_1,\dots,\lambda_\ell)\vdash n$, with parts indexed by positions $[\ell]$. Let $\mu$ be a
coarsening of $\lambda$, and put $m=\ell(\lambda)-\ell(\mu)$. Define $\Lead_{\lambda\mu}(t):=[(s-1)^m]\,c_{\lambda\mu}$.
By Theorem~B (Day~220; proved, and it does not use (N)),
\[
\Lead_{\lambda\mu}=\sum_{\text{tight histories }H\text{ ending at }\mu}\ \prod_{\text{merge steps}} W_k(J).
\]
A history processes positions $\ell,\ell-1,\dots,1$ in turn. At position $i$, where $k=\lambda_i$, it either opens a new
block or merges with $p\ge1$ existing blocks of sizes $J=(j_1,\dots,j_p)$, producing one block of size $k+|J|$. By
Theorem~W (Day~220 \S5b),
\[
W_k(J)=(-1)^p\,[k+|J|]_t\,\frac{\prod_{j\in J}[k]_{t^j}}{[k]_t}.
\]
Opening a block has weight $1$, which is also $W_k(\emptyset)$.

\section*{1. Lemma 1}
\begin{lemma1}
Histories correspond bijectively to increasing forests on $[\ell]$, and the final blocks are the trees.
\end{lemma1}
\begin{proof}
When position $i$ merges with existing blocks, make $i$ the parent of each of those blocks' current roots. Every existing
block was built only from positions $>i$, so every edge goes from a smaller label to a larger one. A history therefore
gives an increasing forest on $[\ell]$, and its trees are the final blocks. Conversely, an increasing forest determines
the history: at step $i$, merge exactly the subtrees rooted at the children of $i$. Those subtrees are complete by then,
because all their labels are $>i$. So the correspondence is a bijection. Histories ending in a single block are
increasing trees, and there are $(\ell-1)!$ of them. A tight history needs $\sum p=m$, which is automatic: each edge is
one unit of $p$, and a forest with $\ell(\mu)$ trees has $\ell-\ell(\mu)$ edges.
\end{proof}

\section*{2. Lemma 2}
\begin{lemma2}
For a vertex $v$, let $N_v$ be the $\lambda$-mass of its subtree. For an increasing tree with root mass $n$,
\[
\prod_v W=\frac{1-t^n}{\prod_i(1-t^{\lambda_i})}\prod_{\text{edges }v\to c}\bigl(t^{\lambda_vN_c}-1\bigr).
\]
\end{lemma2}
\begin{proof}
The merge at $v$ has $k=\lambda_v$ and $J=(N_c)_{c\text{ child of }v}$, so $k+|J|=N_v$. Then
\[
\prod_v W=(-1)^{\#\text{edges}}\prod_v\frac{[N_v]_t}{[\lambda_v]_t}\cdot\prod_{\text{edges }v\to c}[\lambda_v]_{t^{N_c}}.
\]
Here $[N_v]_t/[\lambda_v]_t=(1-t^{N_v})/(1-t^{\lambda_v})$ and
$[\lambda_v]_{t^{N_c}}=(1-t^{\lambda_vN_c})/(1-t^{N_c})$. Every non-root $c$ appears exactly once as a child, so its
factor $(1-t^{N_c})$ cancels, leaving only the root's $(1-t^n)$. The claim follows from
$(-1)^{\#\text{edges}}\prod(1-x)=\prod(x-1)$. Leaves have $N_v=\lambda_v$, which gives the factor $1$, consistent with
opening a block.
\end{proof}

\section*{3. Lemma 3}
\begin{lemma3}
For indeterminates $x_{ij}=x_{ji}$ ($i\ne j$ in $[\ell]$),
\[
\sum_{\text{increasing trees }T\text{ on }[\ell]}\ \prod_{\text{edges }v\to c}\Bigl(\prod_{j\in T_c}x_{vj}-1\Bigr)
=\sum_{\text{connected graphs }G\text{ on }[\ell]}\ \prod_{e\in G}(x_e-1),
\]
where $T_c$ is the vertex set of the subtree at $c$.
\end{lemma3}
\begin{proof}
Map a connected graph $G$ to a tree recursively. The root is $\min V=1$. The components $C_1,\dots,C_r$ of $G-1$ become
the child subtrees, with their minima as children, and we recurse on each $G[C_a]$. Each $C_a$ is connected, so the
recursion is defined. Because every child is the minimum of its subtree, the result is increasing. Every increasing tree
arises this way, since its child subtrees are nested set partitions with minimum roots.

Fix $T$ and sum over the graphs $G$ that map to it. At each vertex $v$, the edges from $v$ into the child set $T_c$ form
an arbitrary \emph{nonempty} subset $S\subseteq T_c$. It must be nonempty because $T_c$ has to be attached to $v$. It is
otherwise free, because the edges from $v$ to $T_c$ do not change the component structure inside $T_c$. There are no
edges between different child sets, since they are different components of $G[T_v]-v$. The edges inside $T_c$ are
accounted for by the recursion. Hence the fibre sum factors, and for each edge $v\to c$
\[
\sum_{\emptyset\ne S\subseteq T_c}\ \prod_{j\in S}(x_{vj}-1)=\prod_{j\in T_c}x_{vj}-1.\qedhere
\]
\end{proof}
\noindent\emph{Specialization.} Setting $x_{ij}=t^{\lambda_i\lambda_j}$ gives $\prod_{j\in T_c}x_{vj}=t^{\lambda_vN_c}$.

\section*{4. Theorem G}
\begin{thmG}
For $\lambda\vdash n$ with $\ell\ge2$,
\[
\Lead_{\lambda,(n)}(t)=\frac{1-t^n}{\prod_i(1-t^{\lambda_i})}\,K_\lambda(t),\qquad
K_\lambda:=\sum_{G\text{ connected on }[\ell]}\ \prod_{ij\in G}\bigl(t^{\lambda_i\lambda_j}-1\bigr).
\]
\end{thmG}
\begin{proof}
Apply Lemma~1 restricted to trees, then Lemma~2, then Lemma~3 with the specialization.
\end{proof}
By the exponential formula, $K_\lambda$ is the connected part (joint cumulant) of
$t^{e_2(\lambda)}=\prod_{i<j}\bigl(1+(t^{\lambda_i\lambda_j}-1)\bigr)$, i.e.\ the logarithm of a group-like element in
the graph Hopf algebra.

\paragraph{Corollaries.}
\begin{itemize}
\item $\lambda=1^n$: the prefactor is $[n]_t/(1-t)^{n-1}$, and
$K=\sum_{\text{conn}}(t-1)^{|E|}=(t-1)^{n-1}I_n(t)$, where $I_n$ is the Mallows--Riordan tree inversion enumerator.
So $\Lead=(-1)^{n-1}[n]_t\,I_n(t)$.
\item $t=0$: every edge factor is $-1$, so $\sum_{\text{conn}}(-1)^{|E|}=(-1)^{\ell-1}(\ell-1)!$. Alternatively, count
the trees directly.
\item $t\to1$: with $\varepsilon=1-t$, the tree form gives
$\Lead\to(-1)^{\ell-1}\,\frac{n}{\prod\lambda_i}\sum_T\prod_{\text{edges }v\to c}\lambda_vN_c$. The claimed value
$(-1)^{\ell-1}n^{\ell-1}$ is equivalent to $\sum_T\prod\lambda_vN_c=n^{\ell-2}\prod\lambda_i$, a weighted Cayley
identity. It checks at $\ell=2$ (giving $\lambda_1\lambda_2$). The general case is a claim from the PROVE session that
is \emph{not re-derived here}; please check it before citing.
\item Sign: for $0<t<1$ the prefactor is positive and each of the $\ell-1$ edge factors is negative. For $t>1$ the
prefactor has sign $(-1)^{\ell-1}$ and the edge factors are positive. Either way
$\operatorname{sgn}\Lead=(-1)^{\ell-1}$ for all $t>0$, $t\ne1$.
\end{itemize}

\section*{5. Theorem F}
\begin{thmF}
For $\mu$ a coarsening of $\lambda$,
\[
\Lead_{\lambda\mu}=\sum_{\substack{\text{set partitions }\pi\text{ of }[\ell]\\\text{with sorted block sums}=\mu}}\ \prod_{B\in\pi}\Lead_{\lambda_B,(|\lambda_B|)},
\]
where a singleton block contributes $1$. There is no interleaving factor.
\end{thmF}
\begin{proof}
By Lemma~1, increasing forests correspond to a set partition together with an increasing tree on each block. The
weight is a product over the trees, and by Lemma~2 each tree's weight depends only on its own parts. Sum.
\end{proof}

\section*{6. Theorem W sober re-check --- lost}
The original \S6 was a line-by-line re-derivation of Day~220 \S5b: signs, top-symbol count, Macdonald III.7 Ex.~2 at
$n=2$, III.2 Ex.~1 multiplicities, and Cauchy. It was lost in the truncation. The session log records ``no gap'', but the
text no longer exists, so the re-check should be treated as unrecorded.

\section*{7. Computation (scripts survive)}
\begin{itemize}
\item \file{scripts/day221/tree_identity.py} $\to$ \file{tree_identity.log}: Lemma~3 with random rational generic
$x_{ij}$, 15/15 ($\ell=2..6$). The history sum (Theorem~W weights) equals the Theorem~G closed form in exact symbolic $t$
(SymPy), \textbf{37/37 for all $\lambda\vdash n\le7$ with $\ell\ge2$}.
\item \file{scripts/day221/tree_identity_n8.py} $\to$ \file{tree_identity_n8.log}: history $=$ telescoped tree form
(Lemma~2) $=$ Theorem~G, in exact rationals at $t\in\{3/5,7/3,-2,2\}$, for all $\lambda\vdash n\le8$ with
$2\le\ell\le6$: \textbf{220/220}.
\item Earlier wake data: G 58/58 ($n\le8$); F 123/123 ($n\le7$); subset-engine agreement 32/32 and 132/132.
\end{itemize}

\section*{8. Gaps}
\begin{itemize}
\item G and F have no gaps beyond their premises: Theorem~B (proved, does not use (N)), Theorem~W (proved; the re-check
text was lost, see \S6) and (KF) (proved, Day~216b).
\item The $t=1$ corollary value is not re-derived here.
\item Novelty has not been searched. Lemma~3 looks classical (a tree--graph / Penrose-type decomposition; Gessel--Wang /
Mallows--Riordan at uniform weights), so the new content would be identifying the leading coefficient of $\star$ with
it. Search targets: ``connected graphs'' with weight $t^{ab}-1$ and the prefactor $(1-t^n)/\prod(1-t^{\lambda_i})$;
coloured inversion enumerators; $q$-Cayley.
\end{itemize}

\end{document}
